\section*{Abstract}
\addcontentsline{toc}{section}{Abstract}

Security operations teams rely on structured threat intelligence
(IOCs, MITRE ATT\&CK mappings, STIX/MISP-ready artifacts) to act on
malware findings quickly. Broad reputation aggregators such as
VirusTotal deliver that structure at enormous scale but do not
reconcile their many signals against each other, so producing a
defensible, evidence-linked technique mapping for one specific sample
usually still falls to a human reverse engineer tracing code and
writing a report by hand. This thesis presents \textbf{MalWhere}, a
modular pipeline that combines static reverse engineering and dynamic
sandbox detonation, reconciles their findings into confidence-scored
MITRE ATT\&CK technique mappings, and exports the result as STIX~2.1
bundles and live MISP events.

The pipeline is validated against three malware families with
distinct architectures: a ransomware binary (Akira); a multi-stage
loader that drops a rootkit driver, an AV-killer, and a remote-access
C2 client (RoningLoader); and a feature-rich .NET infostealer
(WhiteSnakeStealer). Each is ground-truthed through function-level,
Ghidra-verified manual reverse engineering rather than superficial
scanning. Two methodological contributions extend beyond a
conventional static+dynamic pipeline: a cross-source \emph{and}
cross-level confidence reconciliation model that treats a technique
observed by both static and dynamic evidence, or by a parent and a
sub-technique from different sources, as stronger signal than either
alone; and a \emph{resubmission loop} that independently analyzes
every component a multi-stage sample drops or injects, without
misattributing a dropped payload's own capabilities to the parent
binary's confidence score.

The evaluation methodology goes beyond raw precision and recall: every
false positive found during validation was individually traced to its
exact evidence (a specific import combination, a sandbox signature's
raw data field, a string pattern) and checked against the full text of
the corresponding manual report, not just its summary table. Of 15
false positives found this way across the three samples, twelve were
genuine pipeline errors with identifiable, fixed root causes (ten
outright, plus two confidence-calibration corrections whose underlying
generic evidence remains present, honestly downgraded, rather than
disappearing), one was a gap in how ground truth was extracted rather
than a pipeline error, and two were deliberately left open pending
independent confirmation rather than resolved on the pipeline's own
evidence alone. Static detection coverage was subsequently extended
from roughly 30 to approximately 85 MITRE ATT\&CK technique IDs, every
addition sourced from MITRE's own technique descriptions and verified
not to introduce new findings on the validated samples before being
trusted.

The resulting pipeline reaches family-level F1 scores of 0.93 (Akira),
0.81 (WhiteSnakeStealer), and 0.89 (RoningLoader, counting its
resubmitted components), with high precision across all three
samples. The thesis argues that this audit trail, every rule tied to
verifiable evidence, every error diagnosed to a specific cause, every
extension checked for regression, is what makes the pipeline's output
usable for exactly the decision a reputation aggregator's unreconciled
signals cannot support on their own, and is the property most likely
to transfer to malware families the pipeline has not yet seen.

\bigskip
\noindent\textbf{Keywords:} malware analysis, threat intelligence,
MITRE ATT\&CK, STIX, MISP, static analysis, dynamic analysis, sandbox
detonation, reverse engineering, confidence reconciliation.
